\chapter{Kinematic Shear, Tensor Mechanics, and Field Unification}
\label{app:kinematic_shear}

\section*{Abstract}

This appendix develops the macroscopic gravitational-wave sector of the active State-I vacuum and relates it to the microscopic FCC scaffold mechanics established in \hyperref[app:fcc_microphysics]{Appendix~\ref*{app:fcc_microphysics}} and the coherent condensate kinematics developed in \hyperref[app:condensate_impedance]{Appendix~\ref*{app:condensate_impedance}}. The microscopic FCC scaffold possesses a finite transverse restoring coefficient and a primitive-cell inertia normalization whose ratio fixes the invariant unstrained propagation capacity $c_{\mathrm{E}}$. The interstitial condensate is carried with the scaffold without relative flow in the conservative coherent branch and supplies the corresponding volume constraint, but no independently additive condensate inertia is assumed in the microscopic scaffold normalization.

A representative GW150914 calibration is used here to define an effective macroscopic gravitational-wave response and impedance scale. These quantities are not interpreted as direct measurements of intrinsic microscopic vacuum properties or as established universal material constants. They are conditional continuum-response parameters obtained under the adopted gravitational-wave strain--energy-flux mapping.

The dimensionless quantity $\zeta_{\mathrm{GW}}$ is introduced in this appendix only as a diagnostic comparison between the effective macroscopic response coefficient and the baseline microscopic FCC scaffold tangent coefficient. It is not identified with gravitational loading, dynamic wave amplitude, consumed tension capacity, or a derived microscopic participation fraction.

To describe the observed tensor polarizations without identifying an ordinary vector elastic phonon with a spin-2 field, the gravitational wave is represented by a collective transverse-traceless response tensor $\mathsf{H}_{ij}^{(\mathrm{GW})}$, distinct in general from the ordinary Cauchy strain of the FCC displacement field. Its $(+)$ and $(\times)$ components obey the spin-weight-2 transformation law. Static gravity and gravitational radiation are thereby represented as equilibrium and time-dependent sectors of the same State-I constitutive architecture, while the microscopic derivation of the collective quadrupolar projection remains an explicit theory-completion problem.

\section{Microscopic Scaffold Baseline and Macroscopic Gravitational-Wave Response}

\hyperref[app:fcc_microphysics]{Appendix~\ref*{app:fcc_microphysics}} establishes the gravity-independent microscopic FCC scaffold carrier of transverse propagation. Under the leading isotropic contact closure,
\begin{equation}
    \kappa_{\mathrm{n}}=\kappa_{\mathrm{t}},
\end{equation}
the primitive-cell scaffold coefficients satisfy
\begin{equation}
    \rho_{\mathrm{sh}}^{(\mathrm{micro},0)} = \beta_{\mathrm{geo}}\rho_{\mathrm{D}},
\end{equation}
and
\begin{equation}
    K_{\mathrm{sh}}^{(\mathrm{micro},0)} = \rho_{\mathrm{sh}}^{(\mathrm{micro},0)}c_{\mathrm{E}}^2 = \beta_{\mathrm{geo}}\rho_{\mathrm{D}}c_{\mathrm{E}}^2.
\end{equation}

Here $\rho_{\mathrm{D}}$ is the density assigned to the solid nodal phase, whereas $\rho_{\mathrm{sh}}^{(\mathrm{micro},0)}$ is the corresponding nodal mass averaged over the complete FCC primitive-cell volume. The coefficient $K_{\mathrm{sh}}^{(\mathrm{micro},0)}$ is the unstrained primitive-cell tangent shear modulus of the FCC scaffold. It is a constitutive stiffness and is not itself a static shear stress, a gravitational-wave stress, or a measure of gravitational loading.

The quantities above belong to the microscopic scaffold description. A macroscopic gravitational-wave response is described at a different level and is represented here by an effective tensor-response coefficient
\begin{equation}
    K_{\mathrm{GW}}^{(\mathrm{eff})}.
\end{equation}

When an equivalent macroscopic inertia coefficient is useful, it is defined through the unstrained wave-speed parametrization
\begin{equation}
    \rho_{\mathrm{GW}}^{(\mathrm{eff})} \equiv \frac{ K_{\mathrm{GW}}^{(\mathrm{eff})} }{ c_{\mathrm{E}}^2 }.
\end{equation}
Thus,
\begin{equation}
    \frac{ K_{\mathrm{GW}}^{(\mathrm{eff})} }{ \rho_{\mathrm{GW}}^{(\mathrm{eff})} } = c_{\mathrm{E}}^2.
\end{equation}

The quantity $\rho_{\mathrm{GW}}^{(\mathrm{eff})}$ is therefore an equivalent continuum inertia parameter associated with the adopted macroscopic response normalization. It is not interpreted as a literal equilibrium rest-mass density of the State-I vacuum.

Gravitational preloading is described independently by the structural coupling
\begin{equation}
    \alpha_{\mathrm{G}}(\mathbf{x}),
\end{equation}
which measures the local static gravitational loading of the active Soft Floor. In a gravitationally loaded State-I region, the endpoint-transfer law developed in \hyperref[app:gravitational_endpoint_transfer]{Appendix~\ref*{app:gravitational_endpoint_transfer}} gives
\begin{equation}
    c_{\mathrm{P}}(\mathbf{x}) = c_{\mathrm{E}} \left[ 1-\alpha_{\mathrm{G}}(\mathbf{x}) \right]^2.
\end{equation}

The structural coupling $\alpha_{\mathrm{G}}$ and the diagnostic micro-to-macro ratio introduced later in this appendix therefore represent different physical concepts. The former is a local environmental loading variable; the latter compares response coefficients defined at different descriptive scales.

\section{The Coherent State-I Wave Sector}

\hyperref[app:condensate_impedance]{Appendix~\ref*{app:condensate_impedance}} establishes the conservative coherent kinematics of the dual-phase State-I vacuum. Let $\mathbf{u}$ denote the FCC scaffold displacement and $\mathbf{v}_{\mathrm{IC}}$ the interstitial-condensate velocity. Their relative velocity is
\begin{equation}
    \mathbf{v}_{\mathrm{rel}} = \mathbf{v}_{\mathrm{IC}} - \dot{\mathbf{u}}.
\end{equation}

For the coherent branch,
\begin{equation}
    \mathbf{v}_{\mathrm{rel}} = \mathbf{0}.
\end{equation}

Thus the condensate is not assumed to remain motionless. It is carried with the FCC scaffold at the retained conservative order, so coherent propagation occurs without relative condensate flow.

The same selected coherent branch satisfies the incremental volume constraint
\begin{equation}
    \nabla\cdot\mathbf{u} = 0.
\end{equation}

Within the translational acoustic sector, this excludes an ordinary finite-frequency volumetric compression branch and leaves two transverse translational polarizations,
\begin{equation}
    \mathbf{k}\cdot\mathbf{e}_T = 0.
\end{equation}

This statement applies only to the selected translational coherent sector. Rotational, torsional, phase, chiral, and localized relative-flow branches are not thereby removed from the complete microscopic spectrum.

The transverse translational carrier does not, by itself, establish a spin-2 gravitational wave. Two transverse vector displacement modes and a rank-two quadrupolar tensor mode are mathematically distinct objects. Their relation requires a separate collective projection.

\section{GW150914 as a Representative Macroscopic Calibration}

GW150914 provides a representative observational scale for the macroscopic gravitational-wave response. The detected signal swept through approximately $35$--$250\ \mathrm{Hz}$ and reached a strain amplitude of order
\begin{equation}
    h_{\mathrm{GW}} \sim 10^{-21}.
\end{equation}

The source distance and radiated gravitational-wave energy are obtained through waveform reconstruction rather than by a direct local measurement of mechanical energy flux at the detector. For a representative event-scale calibration, take
\begin{equation}
    h_{\mathrm{GW}} \sim 10^{-21},
\end{equation}
and a characteristic frequency scale
\begin{equation}
    f \sim 150\ \mathrm{Hz}, \qquad \omega \sim 9.4\times10^2\ \mathrm{s^{-1}}.
\end{equation}

Using representative radiated-energy, distance, and signal-duration scales gives an isotropic-equivalent flux of order
\begin{equation}
    I_{\mathrm{eff}} \sim 10^{-3}\ \mathrm{W\,m^{-2}}.
\end{equation}

This quantity is used only as a source-model-dependent response calibration. It is not an independently measured material-energy flux through the State-I vacuum.

The distinction is essential. The interferometric observable is gravitational-wave strain, whereas radiated energy, source distance, polarization content, and corresponding energy flux depend on waveform reconstruction. The calculation below therefore defines an effective continuum response scale conditional on the adopted strain--flux mapping; it does not directly measure a universal shear modulus of the vacuum.

\section{The Collective Gravitational-Wave Response Tensor}

The ordinary infinitesimal Cauchy strain associated with an FCC displacement field is
\begin{equation}
    \varepsilon_{ij}^{(u)} = \frac{1}{2} \left( \partial_i u_j + \partial_j u_i \right).
\end{equation}

For an ordinary transverse vector wave traveling along $z$, this strain generally contains components such as $\varepsilon_{xz}^{(u)}$ or $\varepsilon_{yz}^{(u)}$. It does not automatically generate the $(+)$ and $(\times)$ quadrupolar patterns of a gravitational wave.

TFT therefore distinguishes the microscopic translational strain from the collective macroscopic gravitational-wave response tensor
\begin{equation}
    \mathsf{H}_{ij}^{(\mathrm{GW})}.
\end{equation}

The tensor $\mathsf{H}_{ij}^{(\mathrm{GW})}$ is the long-wavelength quadrupolar response variable used to represent the interferometric gravitational-wave deformation. At the present level of the theory, it is a coarse-grained collective field and is not assumed to be identical, component by component, to $\varepsilon_{ij}^{(u)}$.

For propagation along the $z$-axis, define
\begin{equation}
    \mathbf{E}_{+} = 
    \begin{pmatrix}
        1 & 0 & 0 \\
        0 & -1 & 0 \\
        0 & 0 & 0
    \end{pmatrix},
\end{equation}
and
\begin{equation}
    \mathbf{E}_{\times} = 
    \begin{pmatrix}
        0 & 1 & 0 \\
        1 & 0 & 0 \\
        0 & 0 & 0
    \end{pmatrix}.
\end{equation}

The collective response is written
\begin{equation}
    \mathsf{H}_{ij}^{(\mathrm{GW})} = \frac{1}{2} \left[ h_{+}(\mathbf{E}_{+})_{ij} + h_{\times}(\mathbf{E}_{\times})_{ij} \right].
\end{equation}

For a pure $(+)$ mode,
\begin{equation}
    \mathsf{H}_{xx}^{(\mathrm{GW})} = +\frac{h_{+}}{2}, \qquad \mathsf{H}_{yy}^{(\mathrm{GW})} = -\frac{h_{+}}{2},
\end{equation}
so that
\begin{equation}
    \mathsf{H}_{xx}^{(\mathrm{GW})} - \mathsf{H}_{yy}^{(\mathrm{GW})} = h_{+}.
\end{equation}

This fixes the normalization between the collective response tensor and the differential interferometric strain amplitude.

By construction,
\begin{equation}
    \mathsf{H}_{ij}^{(\mathrm{GW})} = \mathsf{H}_{ji}^{(\mathrm{GW})},
\end{equation}
\begin{equation}
    s^i\mathsf{H}_{ij}^{(\mathrm{GW})} = 0,
\end{equation}
and
\begin{equation}
    \mathsf{H}_{ii}^{(\mathrm{GW})} = 0.
\end{equation}

The unresolved microscopic problem is the derivation of a definite projection
\begin{equation}
    \left\{ \text{microscopic FCC contact, translational, rotational, and phase fields} \right\} \longrightarrow \mathsf{H}_{ij}^{(\mathrm{GW})}.
\end{equation}

That projection must ultimately emerge from the collective FCC eigenstructure rather than being imposed on a single primitive-cell translational mode.

\section{Effective Macroscopic Gravitational-Wave Response and Impedance}

At the macroscopic tensor-response level, introduce the effective work-conjugate closure
\begin{equation}
    \delta\sigma_{ij}^{(\mathrm{GW})} = 2K_{\mathrm{GW}}^{(\mathrm{eff})} \mathsf{H}_{ij}^{(\mathrm{GW})}.
\end{equation}

This equation defines an effective long-wavelength tensor-response coefficient for the selected macroscopic branch. It does not assert that $K_{\mathrm{GW}}^{(\mathrm{eff})}$ is the microscopic FCC shear modulus or that the ordinary Cauchy shear mode is itself a spin-2 excitation.

Within this quadratic effective closure,
\begin{equation}
    u_{\mathrm{el}}^{(\mathrm{GW})} = \frac{1}{2} \delta\sigma_{ij}^{(\mathrm{GW})} \mathsf{H}_{ij}^{(\mathrm{GW})}.
\end{equation}

Since
\begin{equation}
    \mathsf{H}_{ij}^{(\mathrm{GW})} \mathsf{H}_{ij}^{(\mathrm{GW})} = \frac{1}{2} \left( h_{+}^2+h_{\times}^2 \right),
\end{equation}
the corresponding elastic contribution is
\begin{equation}
    u_{\mathrm{el}}^{(\mathrm{GW})} = \frac{ K_{\mathrm{GW}}^{(\mathrm{eff})} }{2} \left( h_{+}^2+h_{\times}^2 \right).
\end{equation}

For the adopted traveling-wave normalization, define
\begin{equation}
    h_{\mathrm{eff}}^2 \equiv h_{+,0}^2 + h_{\times,0}^2,
\end{equation}
and write the mean effective flux as
\begin{equation}
    \left\langle I_{\mathrm{GW}}\right\rangle \equiv \frac{ K_{\mathrm{GW}}^{(\mathrm{eff})}c_{\mathrm{E}} }{2} h_{\mathrm{eff}}^2.
\end{equation}

This relation is used operationally to normalize the equivalent macroscopic response coefficient. Therefore,
\begin{equation}
    K_{\mathrm{GW}}^{(\mathrm{eff})} = \frac{ 2I_{\mathrm{eff}} }{ c_{\mathrm{E}}h_{\mathrm{eff}}^2 }.
\end{equation}

Using the representative GW150914 response calibration adopted here,
\begin{equation}
    I_{\mathrm{eff}} \sim 10^{-3}\ \mathrm{W\,m^{-2}}, \qquad h_{\mathrm{eff}} \sim 10^{-21},
\end{equation}
gives
\begin{equation}
    K_{\mathrm{GW}}^{(\mathrm{eff})} \sim 10^{31}\ \mathrm{Pa}.
\end{equation}

The corresponding effective impedance is
\begin{equation}
    Z_{\mathrm{GW}}^{(\mathrm{eff})} \equiv \frac{ K_{\mathrm{GW}}^{(\mathrm{eff})} }{ c_{\mathrm{E}} },
\end{equation}
giving
\begin{equation}
    Z_{\mathrm{GW}}^{(\mathrm{eff})} \sim 3\times10^{22}\ \mathrm{kg\,m^{-2}\,s^{-1}}.
\end{equation}

The equivalent inertia parameter is
\begin{equation}
    \rho_{\mathrm{GW}}^{(\mathrm{eff})} \equiv \frac{ K_{\mathrm{GW}}^{(\mathrm{eff})} }{ c_{\mathrm{E}}^2 },
\end{equation}
so that
\begin{equation}
    \rho_{\mathrm{GW}}^{(\mathrm{eff})} \sim 10^{14}\ \mathrm{kg\,m^{-3}}.
\end{equation}

The quantities
\begin{equation}
    K_{\mathrm{GW}}^{(\mathrm{eff})}, \qquad Z_{\mathrm{GW}}^{(\mathrm{eff})}, \qquad \rho_{\mathrm{GW}}^{(\mathrm{eff})}
\end{equation}
are not three independent measurements. They are three representations of the same adopted macroscopic response normalization:
\begin{equation}
    K_{\mathrm{GW}}^{(\mathrm{eff})} = Z_{\mathrm{GW}}^{(\mathrm{eff})}c_{\mathrm{E}},
\end{equation}
and
\begin{equation}
    \rho_{\mathrm{GW}}^{(\mathrm{eff})} = \frac{ Z_{\mathrm{GW}}^{(\mathrm{eff})} }{ c_{\mathrm{E}} }.
\end{equation}

\subsection*{Spectral and Constitutive Status of the Calibration}

The adopted flux is obtained from a gravitational-wave source model rather than from an independent local material-energy measurement. Consequently, the inferred response coefficient should generally be regarded as
\begin{equation}
    K_{\mathrm{GW}}^{(\mathrm{eff})} = K_{\mathrm{GW}}^{(\mathrm{eff})}(\omega)
\end{equation}
until its frequency dependence is determined from the microscopic theory or demonstrated observationally to approach a constant long-wavelength limit.

The present calculation therefore establishes a representative effective response scale near the selected GW150914 frequency range. It does not establish a universal State-I shear modulus.

A complete TFT treatment must determine whether
\begin{equation}
    \lim_{kr_{\mathrm{D}}\rightarrow0} K_{\mathrm{GW}}^{(\mathrm{eff})}(\omega)
\end{equation}
approaches a source-independent constitutive coefficient and, if so, how that coefficient emerges from the microscopic FCC dynamics.

\section{Diagnostic Microscopic-to-Macroscopic Comparison}

The baseline primitive-cell scaffold normalization derived in \hyperref[app:fcc_microphysics]{Appendix~\ref*{app:fcc_microphysics}} gives
\begin{equation}
    \rho_{\mathrm{sh}}^{(\mathrm{micro},0)} \sim 10^{62}\ \mathrm{kg\,m^{-3}},
\end{equation}
and
\begin{equation}
    K_{\mathrm{sh}}^{(\mathrm{micro},0)} \sim 10^{79}\ \mathrm{Pa}.
\end{equation}

The representative macroscopic gravitational-wave normalization above gives an effective response coefficient many orders of magnitude smaller. To record this scale difference, define
\begin{equation}
    \zeta_{\mathrm{GW}} \equiv \frac{ K_{\mathrm{GW}}^{(\mathrm{eff})} }{ K_{\mathrm{sh}}^{(\mathrm{micro},0)} }.
\end{equation}

For the representative GW150914 calibration,
\begin{equation}
    \zeta_{\mathrm{GW}} \sim 10^{-48}.
\end{equation}

At the present stage, $\zeta_{\mathrm{GW}}$ is a diagnostic ratio between coefficients defined at two different descriptive scales. It is not a gravitational strain, not a dynamic wave amplitude, not a fraction of the available Soft Floor tension capacity, and not yet a derived microscopic participation coefficient.

In particular,
\begin{equation}
    \zeta_{\mathrm{GW}} \neq \alpha_{\mathrm{G}}.
\end{equation}

The structural coupling $\alpha_{\mathrm{G}}$ is a local environmental loading variable determined by the static gravitational state. By contrast, $\zeta_{\mathrm{GW}}$ presently records only the ratio between an effective macroscopic GW response normalization and the baseline microscopic FCC scaffold tangent coefficient.

If the equivalent macroscopic inertia parameter is defined by
\begin{equation}
    \rho_{\mathrm{GW}}^{(\mathrm{eff})} = \frac{ K_{\mathrm{GW}}^{(\mathrm{eff})} }{ c_{\mathrm{E}}^2 },
\end{equation}
then, because
\begin{equation}
    \rho_{\mathrm{sh}}^{(\mathrm{micro},0)} = \frac{ K_{\mathrm{sh}}^{(\mathrm{micro},0)} }{ c_{\mathrm{E}}^2 },
\end{equation}
one obtains algebraically
\begin{equation}
    \frac{ \rho_{\mathrm{GW}}^{(\mathrm{eff})} }{ \rho_{\mathrm{sh}}^{(\mathrm{micro},0)} } = \frac{ K_{\mathrm{GW}}^{(\mathrm{eff})} }{ K_{\mathrm{sh}}^{(\mathrm{micro},0)} } = \zeta_{\mathrm{GW}}.
\end{equation}

This equality is not an independent microscopic result. It follows from representing both scales using the same unstrained stiffness-to-inertia ratio
\begin{equation}
    \frac{K}{\rho} = c_{\mathrm{E}}^2.
\end{equation}

Likewise,
\begin{equation}
    K_{\mathrm{GW}}^{(\mathrm{eff})} = \zeta_{\mathrm{GW}} K_{\mathrm{sh}}^{(\mathrm{micro},0)}
\end{equation}
is presently a rewriting of the ratio definition, not a derived coarse-graining law.

A physical interpretation of $\zeta_{\mathrm{GW}}$ as a genuine collective-mode projection would require two additional results. First, the macroscopic quantity $K_{\mathrm{GW}}^{(\mathrm{eff})}$ must be shown to represent a source-independent constitutive response of the long-wavelength tensor branch. Second, a microscopic projection from the FCC eigenstructure to $\mathsf{H}_{ij}^{(\mathrm{GW})}$ must reproduce that response coefficient.

Until those conditions are satisfied, $\zeta_{\mathrm{GW}}$ remains diagnostic.

\section{Separation from the Microscopic Saturation Scale}

The equivalent macroscopic response normalization may also be used to define a characteristic tensor-stress scale. For a pure $(+)$ response,
\begin{equation}
    \mathsf{H}_{xx}^{(\mathrm{GW})} = \frac{ h_{\mathrm{GW}} }{2},
\end{equation}
so the adopted effective closure gives
\begin{equation}
    \delta\sigma_{xx}^{(\mathrm{GW})} = K_{\mathrm{GW}}^{(\mathrm{eff})} h_{\mathrm{GW}}.
\end{equation}

Define the corresponding equivalent wave-stress scale by
\begin{equation}
    \sigma_{\mathrm{GW}}^{(\mathrm{eq})} \sim K_{\mathrm{GW}}^{(\mathrm{eff})} h_{\mathrm{GW}}.
\end{equation}

For the representative GW150914 calibration,
\begin{equation}
    K_{\mathrm{GW}}^{(\mathrm{eff})} \sim 10^{31}\ \mathrm{Pa},
\end{equation}
and
\begin{equation}
    h_{\mathrm{GW}} \sim 10^{-21},
\end{equation}
so that
\begin{equation}
    \sigma_{\mathrm{GW}}^{(\mathrm{eq})} \sim 10^{10}\ \mathrm{Pa}.
\end{equation}

The microscopic FCC gravitational saturation scale is
\begin{equation}
    \tau_{\mathrm{D,bulk}} \approx 2.18\times10^{96}\ \mathrm{Pa}.
\end{equation}

Hence,
\begin{equation}
    \frac{ \sigma_{\mathrm{GW}}^{(\mathrm{eq})} }{ \tau_{\mathrm{D,bulk}} } \sim 10^{-86}.
\end{equation}

Under the adopted equivalent-stress calibration, the observed gravitational-wave amplitude is therefore extremely small relative to the microscopic FCC saturation scale.

This is a scale-separation diagnostic, not a proof of complete lattice stability and not a direct measurement of microscopic Cauchy stress. \hyperref[app:fcc_microphysics]{Appendix~\ref*{app:fcc_microphysics}} leaves the full finite-wave-number stability problem open; the completed microscopic theory must demonstrate
\begin{equation}
    \omega_s^2(\mathbf{k}) \geq 0
\end{equation}
for every physical branch throughout the first Brillouin zone.

\section{Tensor Polarizations and the Spin-2 Representation}

The collective response tensor introduced above contains two transverse-traceless basis components,
\begin{equation}
    \mathbf{E}_{+}, \qquad \mathbf{E}_{\times}.
\end{equation}

Its transformation under a rotation of the transverse coordinate axes through an angle $\psi$ is
\begin{equation}
    h_{+}' = h_{+}\cos(2\psi) + h_{\times}\sin(2\psi),
\end{equation}
and
\begin{equation}
    h_{\times}' = -h_{+}\sin(2\psi) + h_{\times}\cos(2\psi).
\end{equation}

Equivalently,
\begin{equation}
    h_{+}+ih_{\times} \longrightarrow e^{-2i\psi} \left( h_{+}+ih_{\times} \right).
\end{equation}

The conjugate combination transforms with the opposite phase. The appearance of $2\psi$ is the characteristic spin-weight-2 transformation of a transverse rank-two field. Thus,
\begin{equation}
    \mathsf{H}_{ij}^{(\mathrm{GW})} \text{ carries the macroscopic spin-2 tensor representation.}
\end{equation}

This result must not be attributed to incompressibility alone. The coherent translational condition
\begin{equation}
    \nabla\cdot\mathbf{u} = 0
\end{equation}
removes the ordinary scalar volumetric compression branch from the selected vector displacement sector. It does not by itself construct the rank-two gravitational-wave field.

The tracelessness and transversality of
\begin{equation}
    \mathsf{H}_{ij}^{(\mathrm{GW})}
\end{equation}
follow from the collective transverse-traceless tensor construction introduced above.

The microscopic interpretation therefore remains deliberately limited. \hyperref[app:fcc_microphysics]{Appendix~\ref*{app:fcc_microphysics}} establishes the FCC transverse translational carrier, and \hyperref[app:condensate_impedance]{Appendix~\ref*{app:condensate_impedance}} establishes the coherent dual-phase constraint. Neither appendix yet derives the coarse-grained map by which correlated FCC contact, translational, rotational, and phase motion forms the macroscopic rank-two field $\mathsf{H}_{ij}^{(\mathrm{GW})}$.

Therefore,
\begin{equation}
    \text{two microscopic transverse carrier modes} \not\equiv \text{a derived microscopic spin-2 mode}.
\end{equation}

The present appendix establishes the macroscopic tensor representation. Deriving that representation from the complete microscopic FCC spectrum remains a theory-completion requirement.

\section{Mechanical Energy and the Continuum Tensor Budget}

The effective stress and energy expressions developed above define a macroscopic gravitational-wave response closure.

For the selected tensor branch,
\begin{equation}
    \delta\sigma_{ij}^{(\mathrm{GW})} = 2K_{\mathrm{GW}}^{(\mathrm{eff})} \mathsf{H}_{ij}^{(\mathrm{GW})},
\end{equation}
and
\begin{equation}
    u_{\mathrm{el}}^{(\mathrm{GW})} = \frac{1}{2} \delta\sigma_{ij}^{(\mathrm{GW})} \mathsf{H}_{ij}^{(\mathrm{GW})}.
\end{equation}

These expressions define the work-conjugate effective tensor sector associated with the adopted macroscopic normalization. They should not be promoted to an independent validation against a standard gravitational-wave energy density when the same source-model energy flux has already been used to determine $K_{\mathrm{GW}}^{(\mathrm{eff})}$. Such a comparison would be a normalization consistency check, not independent evidence for a material vacuum.

The complete continuum treatment is supplied by \hyperref[app:cosmic_energy_continuum_mechanics]{Appendix~\ref*{app:cosmic_energy_continuum_mechanics}}. There, the State-I material balance contains the full energy density, material momentum density, mechanical energy flux, and stress structure, together with its long-wavelength observable projection.

This appendix therefore provides the dynamic tensor-response sector required by that continuum hierarchy without duplicating the full energy--momentum derivation.

\section{Static and Dynamic Gravity as Constitutive Sectors of State I}

Static gravitational loading and gravitational radiation correspond to different dynamical sectors of the same active State-I constitutive architecture.

Let the State-I mechanical stress be decomposed schematically as
\begin{equation}
    \sigma_{ij}^{(\mathrm{av})} = \sigma_{ij}^{(0)} + \delta\sigma_{ij}.
\end{equation}

The equilibrium part contains the background structural tension and the static source-induced stress configuration. Its spatial variation generates the tension potential developed in Chapter~10,
\begin{equation}
    \mathbf{g} = -\nabla\Phi_{\tau}.
\end{equation}

Thus,
\begin{equation}
    \text{static gravity} = \text{equilibrated restorative stress configuration of State I}.
\end{equation}

When the source configuration changes, the State-I vacuum acquires a time-dependent disturbance. The coherent translational carrier remains subject to
\begin{equation}
    \nabla\cdot\mathbf{u} = 0,
\end{equation}
while the observable gravitational-wave deformation is represented at the macroscopic level by
\begin{equation}
    \mathsf{H}_{ij}^{(\mathrm{GW})}.
\end{equation}

The microscopic projection connecting the underlying coherent carrier to this collective tensor field has not yet been completed. Accordingly, the present appendix does not identify the effective tensor stress $\delta\sigma_{ij}^{(\mathrm{GW})}$ directly with the microscopic Cauchy stress associated with $\mathbf{u}$.

At the macroscopic level,
\begin{equation}
    \text{gravitational radiation} = \text{time-dependent transverse tensor response of State I}.
\end{equation}

The unification is constitutive rather than algebraically circular. This appendix does not independently derive
\begin{equation}
    g = \frac{GM}{r^2}
\end{equation}
by inserting an already-established gravitational strain into a macroscopic constitutive law and cancelling the stiffness-to-inertia ratio. The static inverse-square field is obtained independently from conservation of the restorative tension-gradient flux in Chapter~10.

The role of this appendix is to establish that the static equilibrium field and the propagating tensor perturbation can belong consistently to the same underlying State-I mechanical architecture.

Accordingly,
\begin{equation}
    \text{equilibrium restorative sector} + \text{dynamic tensor-response sector} = \text{two dynamical limits of State I}.
\end{equation}

\section{Status, Predictive Requirements, and Microscopic Completion}

The results of this appendix fall into three distinct categories.

\subsection*{Established Constitutive and Kinematic Structure}

The microscopic FCC scaffold transverse branch satisfies
\begin{equation}
    K_{\mathrm{sh}}^{(\mathrm{micro},0)} = \rho_{\mathrm{sh}}^{(\mathrm{micro},0)} c_{\mathrm{E}}^2.
\end{equation}

The selected coherent dual-phase branch satisfies
\begin{equation}
    \mathbf{v}_{\mathrm{rel}} = \mathbf{0},
\end{equation}
and
\begin{equation}
    \nabla\cdot\mathbf{u} = 0.
\end{equation}

The macroscopic gravitational-wave deformation is represented by the symmetric transverse-traceless collective tensor
\begin{equation}
    \mathsf{H}_{ij}^{(\mathrm{GW})},
\end{equation}
whose $(+)$ and $(\times)$ components transform with spin weight $2$.

Static restorative gravity and dynamic gravitational radiation are represented as equilibrium and time-dependent sectors of the same State-I constitutive architecture.

\subsection*{Conditionally Normalized Macroscopic Quantities}

Using the representative GW150914 response calibration adopted here,
\begin{equation}
    K_{\mathrm{GW}}^{(\mathrm{eff})} \sim 10^{31}\ \mathrm{Pa},
\end{equation}
\begin{equation}
    \rho_{\mathrm{GW}}^{(\mathrm{eff})} \sim 10^{14}\ \mathrm{kg\,m^{-3}},
\end{equation}
and
\begin{equation}
    Z_{\mathrm{GW}}^{(\mathrm{eff})} \sim 3\times10^{22}\ \mathrm{kg\,m^{-2}\,s^{-1}}.
\end{equation}

These are conditional effective response parameters generated by the adopted macroscopic strain--flux normalization. They are not direct measurements of intrinsic microscopic vacuum properties and are not yet established as source-independent material constants.

Their comparison with the microscopic FCC scaffold tangent coefficient gives the diagnostic ratio
\begin{equation}
    \zeta_{\mathrm{GW}} \sim 10^{-48}.
\end{equation}

At the present stage, this number records a cross-scale comparison only.

\subsection*{Open Microscopic and Multiscale Derivations}

Three related problems remain.

First, the full FCC Bloch spectrum must be completed and shown to satisfy global stability,
\begin{equation}
    \omega_s^2(\mathbf{k}) \geq 0.
\end{equation}

Second, the theory must derive the collective quadrupolar projection
\begin{equation}
    \left\{ \varepsilon_{ij}^{(u)}, \text{ rotational, phase, and contact modes} \right\} \longrightarrow \mathsf{H}_{ij}^{(\mathrm{GW})}
\end{equation}
from the microscopic FCC eigenstructure.

Third, the macroscopic response coefficient $K_{\mathrm{GW}}^{(\mathrm{eff})}$ must be tested for source and frequency dependence and related, if possible, to the microscopic scaffold coefficient through a derived coarse-graining operator.

Only if that microscopic-to-macroscopic construction yields a definite constitutive projection may
\begin{equation}
    \zeta_{\mathrm{GW}}
\end{equation}
be promoted from a diagnostic ratio to a physical collective-mode projection coefficient.

Until those calculations are completed, TFT provides a consistent macroscopic tensor-response representation but does not claim a first-principles microscopic derivation of the spin-2 collective mode or of the numerical value of $\zeta_{\mathrm{GW}}$.

\section*{Conclusion}

The active FCC scaffold possesses a microscopic transverse response whose unstrained stiffness-to-inertia ratio satisfies
\begin{equation}
    \frac{ K_{\mathrm{sh}}^{(\mathrm{micro},0)} }{ \rho_{\mathrm{sh}}^{(\mathrm{micro},0)} } = c_{\mathrm{E}}^2.
\end{equation}

The interstitial condensate is carried with the scaffold without relative flow in the selected coherent branch and supplies the associated volume constraint, while the microscopic scaffold normalization remains
\begin{equation}
    \rho_{\mathrm{sh}}^{(\mathrm{micro},0)} = \beta_{\mathrm{geo}}\rho_{\mathrm{D}}.
\end{equation}

Using the representative GW150914 response calibration adopted in this appendix gives
\begin{equation}
    K_{\mathrm{GW}}^{(\mathrm{eff})} \sim 10^{31}\ \mathrm{Pa}, \qquad \rho_{\mathrm{GW}}^{(\mathrm{eff})} \sim 10^{14}\ \mathrm{kg\,m^{-3}},
\end{equation}
under the adopted macroscopic strain--flux normalization.

The difference between the effective macroscopic response scale and the microscopic FCC scaffold coefficient is recorded diagnostically by
\begin{equation}
    \zeta_{\mathrm{GW}} \equiv \frac{ K_{\mathrm{GW}}^{(\mathrm{eff})} }{ K_{\mathrm{sh}}^{(\mathrm{micro},0)} } \sim 10^{-48}.
\end{equation}

This ratio is distinct from the structural gravitational coupling $\alpha_{\mathrm{G}}$. The quantity $\alpha_{\mathrm{G}}$ describes local environmental gravitational loading, whereas $\zeta_{\mathrm{GW}}$ presently compares coefficients defined at microscopic and macroscopic descriptive scales.

The observed gravitational-wave deformation is represented by the collective transverse-traceless tensor
\begin{equation}
    \mathsf{H}_{ij}^{(\mathrm{GW})},
\end{equation}
whose $(+)$ and $(\times)$ components carry the macroscopic spin-weight-2 representation. This collective tensor is not assumed to be identical to the ordinary Cauchy strain of a transverse vector phonon.

Under the adopted effective response normalization, the equivalent GW stress scale remains approximately $10^{86}$ below the microscopic FCC saturation scale, placing the representative perturbation deeply within the linear-response regime relative to that saturation boundary.

Within this hierarchy, static gravity and gravitational radiation are mechanically related without being conflated. Static gravity is the equilibrium restorative stress configuration of the State-I vacuum, whereas gravitational radiation is represented by its propagating time-dependent tensor response.

The remaining first-principles task is to derive the collective quadrupolar projection directly from the microscopic FCC dynamics and determine whether the diagnostic ratio $\zeta_{\mathrm{GW}}$ acquires a genuine mode-projection interpretation or remains only a comparison between differently normalized descriptions.
